\begin{center}
    {\LARGE \textbf{2016无机化学}} \\[2ex]
    {\large 时量180分钟 \quad 满分100分}
\end{center}

\vspace{0.5cm}
\noindent\textbf{注意事项:}
\begin{enumerate}[label=\arabic*., parsep=0pt, itemsep=0pt]
    \item 答题前,考生先将自己的姓名、学号、考场号、座位号填写在答题卡上,将条形码准确粘贴在考生信息条形码粘贴区。
    \item 考试结束后,将本试卷与答题纸一并收回。
    \item 回答各个问题时,将答案写在答题纸上,写在本试卷上无效。
\end{enumerate}

\vspace{0.5cm}
\noindent\textbf{一、完成并配平下列反应方程式（15 pts.）}

\begin{enumerate}[label=\arabic*., start=1, itemsep=1.5ex]
    \item 向热水中加入过氧化钠。
    \item \ce{B2O3}与\ce{CuO}共熔。
    \item 将单质\ce{Si}加入到浓\ce{NaOH}溶液中。
    \item \ce{NH4NO3}于$200^\circ\text{C}$时受热分解。
    \item 常温下向浓硫酸中通入适量\ce{H2S}。
\end{enumerate}

\vspace{0.5cm}
\noindent\textbf{二、解释下列实验事实，需给出相应的方程式（25 pts.）}

\begin{enumerate}[label=\arabic*., start=1, itemsep=1.5ex]
    \item 向\ce{Fe^{3+}}溶液中加入\ce{KSCN}生成血红色物质，但再向其中滴加\ce{Sn^{2+}}溶液，则颜色消失。
    \item 向黄色铬酸钾溶液中加入酸，则溶液转变为橙色。
    \item 氢氟酸溶液不可以使用玻璃瓶存储而氢氧化钠溶液不可以使用磨口玻璃瓶存储。
    \item 向淀粉\ce{KI}溶液中加入\ce{NaClO}溶液，溶液变成蓝色；再向溶液中滴加足够的\ce{Na2SO3}溶液后，蓝色消失。
    \item 向\ce{Na3PO4}溶液中加入\ce{AgNO3}溶液后，生成黄色沉淀；而向\ce{NaPO3}溶液中加入\ce{AgNO3}溶液后，生成白色沉淀。
\end{enumerate}

\vspace{0.5cm}
\noindent\textbf{三、回答下列各题（70 pts.）}

\begin{enumerate}[label=\arabic*., start=1, itemsep=1.5ex]
    \item 说明\ce{F2}、\ce{Cl2}、\ce{Br2}、\ce{I2}的存在状态及颜色，并解释颜色成因。
    \item 写出\ce{CO}，\ce{N2}，\ce{O2}，\ce{F2}的分子轨道式及键级，并结合价键理论说明为什么\ce{CO}比\ce{N2}活泼。
    \item 常温下\ce{BH3}通常不能以单体的形式存在，而以二聚的\ce{B2H6}形式存在。而\ce{BF3}则能够以单体的形式存在。
    \item 试比较下列配合物的稳定性，并说明原因。
    \begin{enumerate}[label=(\alph*), itemsep=1ex]
        \item \ce{NiCl4^{2-}} 与 \ce{PtCl4^{2-}}
        \item \ce{Cu(NH3)4^{2+}} 与 \ce{Zn(NH3)4^{2+}}
    \end{enumerate}
    \item 分别利用\ce{PbO2}，\ce{NaBiO3}与\ce{Mn^{2+}}溶液反应，在酸性条件下，产物各为什么？写出相应方程式。
    \item \ce{NaClO}，\ce{SO2}，\ce{(NH4)2S2O8}均可使品红褪色，请分析原因。
    \item 分析\ce{CO3^{2-}}，\ce{O3}，\ce{NO3-}的杂化方式及形成的共轭$\pi$键。
    \item 试说明\ce{Ti}、\ce{Zr}、\ce{Hf}中哪两种元素化学性质更接近。
    \item 分析说明\ce{NH3}，\ce{N2H4}，\ce{NH2OH}，\ce{N3H}的酸碱变化规律。
    \item \ce{O}和\ce{F}的电子亲和能小于\ce{S}和\ce{Cl}。
    \item 请比较\ce{Na2CO3}和\ce{NaHCO3}的热稳定性及\ce{AgF}与\ce{AgI}的溶解性。
    \item 在298K时，测得原电池：
    \[
    (-)\ce{Cu(s)}|\ce{Cu^{2+}}(1.0 \times 10^{-2}\mathrm{mol\cdot dm^{-3}}) \parallel \ce{Fe^{3+}}(\ce{c^\theta})|\ce{Fe^{2+}}(\ce{c^\theta})|\ce{Pt(+)}
    \]
    的电动势为 $E = 0.49\mathrm{V}$ （$E^\ominus(\ce{Fe^{3+}/Fe^{2+}}) = 0.77\mathrm{V}$）。
    \begin{enumerate}[label=(\arabic*), itemsep=1ex]
        \item 求解该原电池的标电动势 $E^\ominus$。
        \item 写出该原电池的电极反应，并计算其平衡常数 $K^\ominus$。
    \end{enumerate}
\end{enumerate}

\vspace{0.5cm}
\noindent\textbf{四、分离、鉴别与制备（20 pts.）}

\begin{enumerate}[label=\arabic*., start=1, itemsep=1.5ex]
    \item 以\ce{SiO2}，\ce{C}，\ce{Cl2}和\ce{H2}为原料制备高纯硅。
    \item \ce{AgNO3}，\ce{Zn(NO3)2}和\ce{Mg(NO3)2}溶液均为无色溶液，如何用最简单的方法来识别，请写出鉴别的反应式。
    \item 有五种沉淀：\ce{Pb(OH)2}、\ce{Ag2CrO4}、\ce{AgCl}、\ce{BaSO4}、\ce{BaCO3}，分别装在没有标签的五个试剂瓶中，如何用最简单的方法来辨认。
    \item 设计方案将某混合溶液中的\ce{Al^{3+}}、\ce{Cr^{3+}}、\ce{Cu^{2+}}、\ce{Zn^{2+}}、\ce{Fe^{3+}}分离。
\end{enumerate}

\vspace{0.5cm}
\noindent\textbf{五、推断（20 pts.）}

\begin{enumerate}[label=\arabic*., start=1, itemsep=1.5ex]
    \item 化合物\ce{A}为可溶于水的白色固体，将\ce{A}加热后可以生成白色固体\ce{B}和刺激性气体\ce{C}。\ce{C}通入水中生成溶液\ce{D}；\ce{D}中通入\ce{H2S}气体后生成淡黄色沉淀\ce{E}，固体\ce{E}加入到沸腾的\ce{Na2SO3}溶液中生成无色溶液\ce{F}，\ce{F}溶液可以溶解\ce{AgBr}粉末。固体\ce{B}溶于热\ce{HCl}溶液可以生成\ce{G}。\ce{G}与适量\ce{NaOH}溶液作用产生沉淀现象。加入过量\ce{NaOH}溶液后沉淀消失，若\ce{G}与\ce{NH4HS}溶液作用，则生成白色沉淀\ce{H}。在空气中灼烧\ce{H}，会变成原来的白色固体\ce{B}和\ce{C}。若用稀\ce{HCl}溶液与化合物\ce{A}作用，则生成溶液\ce{G}和气体\ce{C}。
    试判断各字母所代表的物质，并写出相关的化学反应方程式。
    \item 化合物\ce{A}为白色固体，加热\ce{A}分解为固体\ce{B}和气体混合物\ce{C}，固体\ce{B}溶于\ce{HNO3}，将\ce{C}通过冰盐水冷却管，可得到无色液体\ce{D}和气体\ce{E}。向热的\ce{A}溶液中滴加硫代乙酰氨（\ce{CH3CSNH2}）溶液后生成黑色沉淀\ce{F}，沉淀\ce{F}与适合浓度的\ce{HCl}稀溶液反应可使沉淀变白；\ce{A}的溶液中加入\ce{KI}溶液得金黄色沉淀\ce{G}，\ce{G}溶于热水。\ce{D}加热变为红棕色气体\ce{H}。\ce{E}与\ce{Hg}（汞）加热燃烧可得\ce{I}，\ce{I}与573K下加热后又可分解生成\ce{E}和\ce{Hg}（汞）。
    试判断各字母所代表的物质，并写出相关的化学反应方程式。
\end{enumerate}